The transgression boundary term makes the declared variational principle and gauge generators differentiable.  Endpoint gauge parameters must match at the common boundary.  Under endpoint Dirichlet data or $\Theta|_{\partial M_5}=0$, the surface variation vanishes.  Under admitted asymptotic symmetries the same term becomes the physical surface charge \citep{MoraOleaTroncosoZanelli2006,PaisSalgadoRebolledoVera2023}.

Nontrivial boundaries may support edge modes.  Those modes are extra physical content; they cannot be counted as a mechanism that removes unwanted bulk modes while leaving pure four-dimensional GR.

\section{Strengthened restricted no-go theorem}
\label{sec:nogo}

\begin{theorem}[Restricted pure-connection obstruction]
Consider a five-dimensional Chern--Simons or fixed-reference transgression theory on $M_4\times S^1$, with gauge group locally $\Spin(4,2)$ and a nondegenerate rank-three invariant polynomial.  Assume:
\begin{enumerate}
\item locality, polynomiality and finitely many five-dimensional fields;
\item pure-connection Chern--Simons/transgression bulk dynamics;
\item a regular compactification retaining every connection zero mode;
\item an invertible field map or genuine gauge/symplectic quotient to the daughter;
\item no external compensator datum that explicitly selects gauge-invariant holonomy information;
\item a regular constant-rank parent canonical sector;
\item nonzero Kaluza--Klein modes removed only by a proved consistent truncation or controlled decoupling limit;
\item endpoint connections and the compensator live in fully declared bundle, cobordism and structure-group sectors.
\end{enumerate}
Then the theory cannot reduce to pure four-dimensional Einstein--Cartan/Einstein--Hilbert gravity with cosmological constant and only two local graviton polarizations.
\end{theorem}

\begin{proof}
The complete connection zero mode contains the adjoint scalar $\phiz=A_y$, and fibre integration produces $\int\langle\phiz F\wedge F\rangle$.  Retaining all zero modes retains $\phiz$ and its equation.  The value $vJ_{54}$ has gauge-invariant conjugacy data, witnessed by $\Tr_{\mathbf6}W=4+2\cosh(L_yv)$, so fixing it is not pure BRST gauge.  Its radial equation $\langle\phiz F\wedge F\rangle=0$ survives and is violated by generic Einstein solutions.  Independently, a regular $\Spin(4,2)$ Chern--Simons sector has thirteen local configuration degrees of freedom.  A regular invertible map or gauge quotient cannot delete eleven physical modes without new constraints or a decoupling mechanism.  Moving to the maximally symmetric Chern--Simons root abandons regularity because $\Om$ and the ordinary quadratic graviton operator vanish.  Retaining $n=0,\pm1$ is not BRST closed; exact nonlinear closure generates the full tower, while the first pair adds twenty-six real modes in a regular quadratic sector.  Finally, although the fixed surface has the correct local Einstein--Cartan plus Euler presymplectic pullback, the parent Hamiltonian flow is not tangent to it because of the retained holonomy equation.  Global existence of the compensator further imposes a Wilson-line conjugacy class and a Lorentz structure-group reduction.  These conclusions contradict assumptions 3--8 if the daughter is pure Einstein gravity.  At least one assumption must be abandoned.
\end{proof}

\paragraph{Legitimate escape routes.}
The theorem is deliberately narrow.  Constrained Stelle--West compensators add a norm constraint; Plebanski/BF theories add two-forms and simplicity constraints; boundary gauged Wess--Zumino--Witten sectors retain edge/group fields; holographic induced gravity is effective and nonlocal from the daughter viewpoint; orbifolds change the mode algebra; a dynamical phase transition adds sector-selection dynamics; and a singular AdS phase abandons constant-rank regularity.  These are not counterexamples because each pays by violating an explicit assumption.

\section{Acceptance-gate audit}
\label{sec:gates}

The verification suite passing its algebraic checks is not the same as the candidate passing TP-00.  The code confirms the calculations; the calculations reject the strong parent.  Scores are independent and are never averaged.

\begin{longtable}{@{}p{0.10\textwidth}p{0.12\textwidth}p{0.12\textwidth}p{0.58\textwidth}@{}}
\toprule
Gate & Score & Status & Reason \\
\midrule
\endfirsthead
\toprule
Gate & Score & Status & Reason \\
\midrule
\endhead
ALG-01 & 2 & Pass & Complete parent, zero-mode fields, observables and boundary data are declared. \\
ALG-02 & 0 & Fail strong reduction & Full five-dimensional algebra closes, but $n=0,\pm1$ does not; the fixed Einstein surface is not preserved by BRST/parent flow. \\
ALG-03 & NA & Not applicable & No chiral matter is claimed. \\
ALG-04 & 1 & Conditional & Bundle and holonomy requirements are explicit; the minimal Lorentzian integral level is global-form dependent. \\
ALG-05 & 0 & Fail & Thirteen regular parent modes versus two in GR; first KK pair adds twenty-six. \\
ALG-06 & 0 & Fail strong reduction & Parent canonical algebra is consistent, but fixing holonomy deletes a gauge-invariant equation and is not a symplectic quotient. \\
PERT-01 & 0 & Fail GR-like vacuum & The maximally symmetric Chern--Simons point has no ordinary quadratic kinetic operator. \\
PERT-02 & 0 & Fail GR-like vacuum & Rank-zero principal kinetic form gives no ordinary strongly hyperbolic graviton problem there. \\
PERT-03 & 1 & Conditional & Classical canonical consistency is not a quantum positive-spectrum construction. \\
PERT-04 & 1 & Conditional & Generic characteristic cones remain background dependent; the fixed daughter inherits GR only after external restriction. \\
PERT-05 & 0 & Fail & No radiative/Wilsonian mechanism decouples omitted zero modes or the KK tower. \\
REC-01 & 0 & Fail full parent & The fixed daughter recovers Einstein--Cartan exactly, but the unrestricted parent does not preserve or dynamically approach it. \\
REC-02 & NA & Not applicable & No Standard Model matter sector is claimed. \\
REC-03 & NA & Not applicable & No particle-pole or mixing prediction is claimed. \\
REC-04 & 2 & Local pass & Local Lorentz and diffeomorphism covariance are maintained; the failure is dynamical rather than explicit symmetry breaking. \\
REP-01 & 2 & Pass & Deterministic code, tests, source ledgers, machine-readable outputs and manifest are supplied. \\
\bottomrule
\end{longtable}

The fatal zeros reject the strong spectrum-preserving parent in its claimed domain.  They do not invalidate the exact cohomological descent identities.

\section{Novelty boundary and publication status}
\label{sec:novelty}

Every broad ingredient has substantial precedent:
\begin{itemize}
\item Chern--Weil descent and transgression forms \citep{ChernSimons1974,MoraOleaTroncosoZanelli2006};
\item higher-dimensional Chern--Simons canonical structure \citep{BanadosGarayHenneaux1995,BanadosGarayHenneaux1996,MiskovicTroncosoZanelli2005};
\item MacDowell--Mansouri and constrained compensators \citep{MacDowellMansouri1977,StelleWest1980};
\item odd-to-even $\Phi F^n$ gravity reductions \citep{Chamseddine1990};
\item Wilson-line conjugacy classes and gauge symmetry breaking \citep{Hosotani2005};
\item Hamiltonian boundary differentiability and covariant phase space \citep{ReggeTeitelboim1974,LeeWald1990,IyerWald1994}.
\end{itemize}

The narrowest defensible project contribution is the combined, reproducible distinction between:
\[
\begin{gathered}
\text{exact cohomological descent},\qquad
\text{correct local symplectic pullback},\\
\text{BRST gauge fixing},\qquad
\text{consistent truncation},\qquad
\text{spectrum-preserving reduction}.
\end{gathered}
\]
Specifically, v1.1 supplies:
\begin{enumerate}
\item an explicit regular-sector numerical witness for the epsilon invariant;
\item the orbit/centralizer and Wilson-line proof that fixed holonomy carries physical data;
\item the gauge-invariant radial equation that survives orientation gauge fixing;
